\section*{الفصل \ref{ch:b1:respiration-fermentation} --- التنفس الخلوي والتخمر}

\begin{solution}{exo:b1:respiration-fermentation:1}
\omterm{def:b1:respiration-fermentation:glycolysis}{التحلل السكري}، في \omterm{def:b1:eukaryotic-cell:organelle}{العصارة الخلوية}: 2 بيروفات، و2 \omterm{def:b1:water-small-molecules:atp}{ATP}، و2 NADH. وأكسدة
البيروفات \omterm{def:b1:respiration-fermentation:krebs}{ودورة كريبس}، في \omterm{def:b1:eukaryotic-cell:mitochondrion}{حشوة المتقدرة}: 6 $\mathrm{CO_2}$، و8 NADH،
و2 $\mathrm{FADH_2}$، و2 GTP. والفسفرة المؤكسدة، في الغشاء الداخلي:
ماء من الأكسجين، ونحو 26 \omterm{def:b1:water-small-molecules:atp}{ATP}.
\end{solution}

\begin{solution}{exo:b1:respiration-fermentation:2}
غلوكوز $+ 2\,\mathrm{NAD^+} + 2\,\mathrm{ADP} + 2\,\mathrm{P_i} \to
2$ بيروفات $+ 2\,\mathrm{NADH} + 2\,\mathrm{H^+} + 2\,\mathrm{ATP} +
2\,\mathrm{H_2O}$. \omterm{def:b1:respiration-fermentation:glycolysis}{والفسفرة على مستوى الركيزة}: أن يُنقل فوسفات مباشرة
من وسيط غني بالطاقة (1،3-ثنائي فوسفوغليسيرات، وفوسفوإينول بيروفات)
إلى ADP \omterm{def:b1:enzymes:enzyme}{بإنزيم}، بلا تدرج غشائي ولا أكسجين.
\end{solution}

\begin{solution}{exo:b1:respiration-fermentation:3}
NADH و$\mathrm{CO_2}$ (من الإيزوسترات إلى $\alpha$-كيتوغلوتارات)،
وNADH و$\mathrm{CO_2}$ (من $\alpha$-كيتوغلوتارات إلى سكسينيل-CoA)،
وGTP (من سكسينيل-CoA إلى سكسينات)، و$\mathrm{FADH_2}$ (من سكسينات
إلى فومارات)، وNADH (من مالات إلى أوكسالوأسيتات). و$\mathrm{CO_2}$
الاثنان يأتيان من نزعي الكربوكسيل من الحمضين السداسي والخماسي.
\end{solution}

\begin{solution}{exo:b1:respiration-fermentation:4}
يجدد $\mathrm{NAD^+}$ حتى يستطيع \omterm{def:b1:respiration-fermentation:glycolysis}{التحلل السكري} أن يستمر ويعطي \omterm{def:b1:water-small-molecules:atp}{ATP}
اثنين لكل غلوكوز. وهو لا يؤكسد الكربون، ولا يستخرج الطاقة الباقية
(فالناتج يحتفظ بجلها)، ولا يصنع أكثر من جزء من خمسة عشر من \omterm{def:b1:water-small-molecules:atp}{ATP}
التنفس.
\end{solution}

\begin{solution}{exo:b1:respiration-fermentation:5}
إلى الأكسجين: $-2\times 96\,500\times 1.14 = \qty{-220}{kJ/mol}$، أي
أربعة \omterm{def:b1:water-small-molecules:atp}{ATP} عند \qty{50}{kJ}. وإلى اليوبيكينون:
$-2\times 96\,500\times 0.37 =
\qty{-71}{kJ/mol}$، أي \omterm{def:b1:water-small-molecules:atp}{ATP} واحد.
\end{solution}

\begin{solution}{exo:b1:respiration-fermentation:6}
العضلة: $2 + 2$ (GTP) $+ 8\times 2.5$ (NADH المتقدري) $+ 2\times
1.5$ ($\mathrm{FADH_2}$) $+ 2\times 1.5$ (NADH العصاري عبر \omterm{def:b1:respiration-fermentation:carriers}{FAD}) $=
30$. والكبد: NADH العصاري يعطي $2\times 2.5$: أي 32.
\end{solution}

\begin{solution}{exo:b1:respiration-fermentation:7}
RQ $= 2.2$. والغلوكوز المتنفَّس $x$: $6x = 1.0$، فيكون
$x = \qty{0.167}{mmol}$؛ و$\mathrm{CO_2}$: $6x + 2y = 2.2$، فيكون
$2y = 1.2$ و$y = \qty{0.6}{mmol}$. فالمخمَّر
$0.6/0.767 = \qty{78}{\%}$ من الغلوكوز، والمتنفَّس \qty{22}{\%}.
\end{solution}

\begin{solution}{exo:b1:respiration-fermentation:8}
يرتفع استهلاك الأكسجين (فالسلسلة، متحررة من الضغط الراجع للتدرج،
تجري بأقصى سرعتها)، ويتوقف تركيب \omterm{def:b1:water-small-molecules:atp}{ATP} (فلا تدرج يقود السنتاز)، وتظهر
كل طاقة الإلكترونات حرارةً. فقد كان المرضى يحرقون شحمهم سريعًا
ويموتون بحرارة جسم لا يخفضها شيء.
\end{solution}

\begin{solution}{exo:b1:respiration-fermentation:9}
تنسد السلسلة: فيصير كل حامل مرجعًا، ولا تُضخ بروتونات، ولا يُصنع \omterm{def:b1:water-small-molecules:atp}{ATP}.
ولا يُعاد أكسدة NADH، فتتوقف \omterm{def:b1:respiration-fermentation:krebs}{دورة كريبس} ونازعة هيدروجين البيروفات
لنقص $\mathrm{NAD^+}$. أما \omterm{def:b1:respiration-fermentation:glycolysis}{التحلل السكري}، متحررًا من تثبيط \omterm{def:b1:water-small-molecules:atp}{ATP}،
فيجري سريعًا ولا بد له من تجديد $\mathrm{NAD^+}$ بإرجاع البيروفات إلى
لاكتات، تغمر الدم.
\end{solution}

\begin{solution}{exo:b1:respiration-fermentation:10}
$8\times 10 + 7\times 2.5 + 7\times 1.5 - 2 = 80 + 17.5 + 10.5 - 2 =
106$ \omterm{def:b1:water-small-molecules:atp}{ATP}؛ و$106/16 = 6.6$ لكل كربون. والأكسجين: البالميتات 23
$\mathrm{O_2}$ مقابل 106 \omterm{def:b1:water-small-molecules:atp}{ATP}، أي \qty{0.22}{O_2} لكل \omterm{def:b1:water-small-molecules:atp}{ATP}؛ والغلوكوز 6
مقابل 30، أي \qty{0.20}{O_2} لكل \omterm{def:b1:water-small-molecules:atp}{ATP}. فلواحدة الأكسجين يكون الغلوكوز
أفضل بنسبة \qty{10}{\%}: فالفقمة، وهي محدودة الأكسجين تحت الماء، تفضل
\omterm{def:b1:carbohydrates:monosaccharide}{الكربوهيدرات}؛ ولواحدة الكتلة يكون \omterm{def:b1:lipids:triglyceride}{الشحم} أفضل بالضعف: فالطنان، وهو
محدود الكتلة في الطيران، يهاجر على \omterm{def:b1:lipids:triglyceride}{الشحم}.
\end{solution}

\begin{solution}{exo:b1:respiration-fermentation:11}
\omterm{def:b1:water-small-molecules:atp}{ATP} وحده: $5/3 < \qty{2}{s}$. والفوسفوكرياتين: $25/3 = \qty{8}{s}$.
\omterm{def:b1:respiration-fermentation:fermentation}{والتخمر}: $2\times 3 = \qty{6}{mmol/L}$ من \omterm{def:b1:water-small-molecules:atp}{ATP} في الثانية، وهي تكفي
الطلب، ما دام \omterm{def:b1:carbohydrates:polysaccharide}{الغليكوجين} واحتمال اللاكتات (عشرات الثواني). أما إيصال
الأكسجين فيستغرق دقيقة أو أكثر ليرتفع ويمد بما لا يزيد على نحو
\qty{1}{mmol/L} من \omterm{def:b1:water-small-molecules:atp}{ATP} في الثانية في ليف: فسباق العشر ثوان ينتهي قبل
أن يبدأ التنفس بالمساعدة، ويدفع ثمنه الفوسفوكرياتين \omterm{def:b1:respiration-fermentation:fermentation}{والتخمر}.
\end{solution}

\begin{solution}{exo:b1:respiration-fermentation:12}
التدرج يخزن الطاقة فرقًا في تركيز البروتونات وفي الكمون الكهربائي عبر
غشاء عازل --- وهو بالضبط مكثفة مشحونة مع \omterm{def:b1:cell-unit-of-life:cell}{خلية} تركيز؛ والسلسلة هي
الشاحن، يقودها هبوط الإلكترونات؛ والسنتاز محرك يديره تيار البروتونات،
وهو بالضبط عنفة، تقرن الجريان بدوران ميكانيكي يصنع \omterm{def:b1:water-small-molecules:atp}{ATP}. وترتخي
المشابهة في أن «البطارية» لا تحمل إلا أجزاء من الألف من الثانية من
استهلاك \omterm{def:b1:cell-unit-of-life:cell}{الخلية} ولا بد من شحنها باستمرار، وأن سائر ناقلات الغشاء نفسه
تسحب منها أيضًا، وأن العنفة تستطيع أن تدور بالمقلوب، فتحلمه \omterm{def:b1:water-small-molecules:atp}{ATP} لتضخ
البروتونات حين تفشل السلسلة.
\end{solution}

\begin{solution}{pb:b1:respiration-fermentation:1}
\textbf{1.} \omterm{def:b1:respiration-fermentation:glycolysis}{التحلل السكري}: 2 \omterm{def:b1:water-small-molecules:atp}{ATP}، و2 NADH (عصاري). ونازعة هيدروجين
البيروفات: 2 NADH. وكريبس: 6 NADH، و2 $\mathrm{FADH_2}$، و2 GTP.
\textbf{2.} NADH المتقدري $8\times 10 = 80$؛ و$\mathrm{FADH_2}$
$2\times 6 = 12$؛ وNADH العصاري بوصفه $\mathrm{FADH_2}$
$2\times 6 = 12$: أي 104 بروتونات.
\textbf{3.} $104/4 = 26$ \omterm{def:b1:water-small-molecules:atp}{ATP}؛ والمجموع $26 + 2 + 2 = 30$.
\textbf{4.} 2 \omterm{def:b1:water-small-molecules:atp}{ATP}.
\textbf{5.} 15.
\textbf{6.} مع الهواء $30\times 50 = \qty{1500}{kJ}$: أي \qty{52}{\%}
من 2870. وبدونه: \qty{100}{kJ}، أي \qty{3.5}{\%} من 2870 لكن
\qty{43}{\%} من 235 المحررة --- \omterm{def:b1:respiration-fermentation:fermentation}{فالتخمر} كفء في التقاط ما يحرره؛
لكنه لا يحرر إلا قليلًا.
\textbf{7.} المهوى $30/30 = \qty{1}{mmol}$ من الغلوكوز في الساعة؛
والمحكم $30/2 = \qty{15}{mmol}$.
\textbf{8.} المحكم: \qty{30}{mmol} من الإيثانول و\qty{30}{mmol} من
$\mathrm{CO_2}$ في الساعة. والمهوى: \qty{6}{mmol} من $\mathrm{CO_2}$
مطلقة و\qty{6}{mmol} من $\mathrm{O_2}$ مستهلكة.
\textbf{9.} $2\times 10.5 = \qty{21}{g}$ لكل مول، أي \qty{0.12}{g}
لكل غرام غلوكوز.
\textbf{10.} $0.5\times 180 = \qty{90}{g}$ لكل مول، أي ما يعادل
$90/10.5 =
8.6$ \omterm{def:b1:water-small-molecules:atp}{ATP} مقابل 30 متاحة: فمركب \omterm{def:b1:water-small-molecules:atp}{ATP} في فائض؛ والمردود يحده
الكربون (إذ يُؤكسد نصف الغلوكوز إلى $\mathrm{CO_2}$ لصنع \omterm{def:b1:water-small-molecules:atp}{ATP}، والكتلة
الحيوية تحتاج إلى هياكل كربونية وقدرة مرجعة)، ويُنفق فائض \omterm{def:b1:water-small-molecules:atp}{ATP} على
الصيانة.
\textbf{11.} للحصول على \omterm{def:b1:water-small-molecules:atp}{ATP} نفسه تستهلك المزرعة المحكمة غلوكوزًا أكثر
بخمس عشرة مرة (السؤال 7) وتنمو خُمسًا لكل غرام (السؤال 9): أي استهلاك
سريع للسكر، ونمو قليل، وكحول --- وهي ملاحظة باستور.
والفسفوفركتوكيناز، المثبَّط بمركب \omterm{def:b1:water-small-molecules:atp}{ATP} وبالسترات عند التهوية، يتحرر حين
يتوقف التنفس.
\textbf{12.} $\mathrm{CO_2}$: $44/44 = \qty{1.0}{mmol}$؛
و$\mathrm{O_2}$: $20/32 = \qty{0.625}{mmol}$: فالحاصل $= 1.6$.
\textbf{13.} $\mathrm{O_2}$: $6x = 0.625$، فيكون
$x = \qty{0.104}{mmol}$. و$\mathrm{CO_2}$: $6x + 2y = 1.0$، فيكون
$2y = 0.375$ و$y = \qty{0.1875}{mmol}$.
\textbf{14.} المخمَّر $0.1875/0.292 = \qty{64}{\%}$ من الغلوكوز. \omterm{def:b1:water-small-molecules:atp}{وATP}
من \omterm{def:b1:respiration-fermentation:fermentation}{التخمر} $2y = 0.375$ مقابل $30x = 3.12$ من التنفس: أي \qty{11}{\%}.
\textbf{15.} $16/23 = 0.70$.
\textbf{16.} صائمًا يحرق الحيوان \omterm{def:b1:lipids:triglyceride}{الشحم} (حاصل 0.7)؛ ومطعَمًا يحرق
\omterm{def:b1:carbohydrates:monosaccharide}{كربوهيدرات} الوجبة (1.0).
\textbf{17.} $96\,500\times 0.22 = \qty{21.2}{kJ/mol}$.
\textbf{18.} $4\times 21.2 = \qty{85}{kJ}$ مقابل \omterm{def:b1:water-small-molecules:atp}{ATP} قدره
\qty{50}{kJ}: أي \qty{59}{\%}.
\textbf{19.} $2\times 96\,500\times 1.14 = \qty{220}{kJ}$؛ والبروتونات
العشرة تخزن $10\times 21.2 = \qty{212}{kJ}$: أي \qty{96}{\%} ---
فالسلسلة لا تهدر إلا قليلًا؛ والفواقد عند السنتاز والناقلات.
\textbf{20.} $0.96\times 0.59 = 0.57$؛ و$2.5\times 50/220 = 0.57$: أي
الرقم نفسه بالطريقين.
\textbf{21.} يرتفع استهلاك الأكسجين، ويهبط مردود \omterm{def:b1:water-small-molecules:atp}{ATP} نحو مردود \omterm{def:b1:respiration-fermentation:fermentation}{التخمر}
(2 لكل غلوكوز من \omterm{def:b1:respiration-fermentation:glycolysis}{التحلل السكري}، وPFK متحرر فيرتفع استهلاك الغلوكوز)،
ويرتفع إنتاج الحرارة، وينهار النمو.
\textbf{22.} يتوقف تركيب \omterm{def:b1:water-small-molecules:atp}{ATP}؛ ولا يستطيع التدرج أن يتفرغ، فتتوقف
السلسلة أمامه، ويتوقف استهلاك الأكسجين، وتخمّر \omterm{def:b1:cell-unit-of-life:cell}{الخلية} ما تستطيع. ثم
إضافة فاصل اقتران تدع البروتونات ترشح راجعة: فتستأنف السلسلة واستهلاك
الأكسجين، وما زال بلا \omterm{def:b1:water-small-molecules:atp}{ATP}.
\textbf{23.} NADH لا يضخ إلا عبر المعقدين III وIV: أي $6/4 = 1.5$؛
ولكل غلوكوز $8\times 1.5 + 2\times 1.5 + 2\times 1.5 + 4 = 22$ \omterm{def:b1:water-small-molecules:atp}{ATP}.
\textbf{24.} بلا أكسجين لا تستطيع السلسلة أن تأخذ الإلكترونات، فلا
يُعاد أكسدة NADH ويبقى المخزون مرجعًا؛ ونازعات الهيدروجين الثلاث
المعتمدة على \omterm{def:b1:respiration-fermentation:carriers}{NAD} في الدورة لا تجد $\mathrm{NAD^+}$ فتتوقف الدورة ---
ولهذا يجب على \omterm{def:b1:respiration-fermentation:fermentation}{التخمر} أن يجدد $\mathrm{NAD^+}$ في العصارة ليستمر
\omterm{def:b1:respiration-fermentation:glycolysis}{التحلل السكري}.
\textbf{25.} 30 \omterm{def:b1:water-small-molecules:atp}{ATP} مع الهواء، و2 بدونه: أي نسبة 15؛ \omterm{thm:b1:respiration-fermentation:chemiosmosis}{ونسبة P/O} تساوي
2.5 عند NADH (10 بروتونات على 4 لكل \omterm{def:b1:water-small-molecules:atp}{ATP}) و1.5 عند $\mathrm{FADH_2}$
(6 على 4).
\end{solution}
